% TEMPLATE-MILC-v3.tex - plantilla maestra UNICA de las guias MILC v3.
%
% La usan las tres familias de guias, cada una con su driver:
%   · Grados 6-11  -> scripts/build-guias-g11.py
%   · Bebras       -> scripts/build-guias-bebras.py
%   · Semillero    -> scripts/build-guias-semillero.py
%
% Antes habia tres copias casi identicas (~400 lineas iguales, 6 distintas):
% cada mejora habia que aplicarla tres veces, y en la practica no se hacia
% (el motor de recursos vivia solo en dos de ellas). Lo que varia entre
% familias esta parametrizado en 6 placeholders que cada driver llena
% (se nombran aqui SIN su sintaxis real: el builder reemplaza por texto plano,
% tambien dentro de los comentarios, y un valor multilinea rompe el preambulo):
%
%   HEADER_IZQ        encabezado izquierdo de pagina
%   HEADER_DER        encabezado derecho de pagina
%   PORTADA_ETIQUETA  rotulo sobre el numero grande (GRADO/MOMENTO/MODULO)
%   PORTADA_SERIE     linea de serie de la portada
%   PORTADA_ID        identificador de la guia en la portada
%   PORTADA_CIRCULO   el circulito de grado (°), o vacio si no aplica
%
% El resto de claves las documenta content/guias/_SCHEMA.md. El script valida
% que no queden placeholders sin reemplazar antes de escribir el archivo final.

\documentclass[11pt,letterpaper]{article}
\usepackage[margin=2.0cm]{geometry}
\usepackage[table]{xcolor}
\usepackage{fontspec}
\usepackage[spanish]{babel}
\usepackage{tikz}
% Librerías TikZ para los diagramas de las guías (bloque `recursos:`):
%  · babel  → babel en español vuelve activos caracteres como «>», y TikZ los usa
%             en sus opciones (>=stealth, ->). Sin esta librería, cualquier
%             diagrama con flechas revienta con «\language@active@arg> has an
%             extra }». Debe ir ANTES de que se dibuje nada.
%  · positioning        → sintaxis «below left=2pt and 3pt».
%  · decorations.pathreplacing → llaves ({brace}) para señalar rangos.
\usetikzlibrary{babel,positioning,decorations.pathreplacing}
\usepackage{graphicx}
% Circuitos: el trabajo de electrónica se hace hoy con micro:bit y sus sensores
% integrados (MakeCode, sin cableado), así que no se necesitan esquemáticos.
% Si algún día entran montajes con protoboard, basta descomentar esta línea:
% los diagramas .tex/.tikz declarados en `recursos:` ya se incrustan solos.
% \usepackage{circuitikz}
\usepackage[most]{tcolorbox}
\usepackage{tabularx}
\usepackage{array}
\usepackage{fancyhdr}
\usepackage{titlesec}
\usepackage[document]{ragged2e}  % bandera a la derecha en todo el cuerpo
\usepackage{enumitem}
\usepackage{microtype}

% Helvetica Neue no trae la flecha U+2192, asi que cada «→» escrito literalmente
% en el YAML se compilaba a NADA, en silencio (462 flechas en 61 guias: rutas del
% tipo «paleta → tipografia → lockup» salian rotas). Mapeamos el caracter a su
% version matematica, que si tiene glifo. Asi el autor puede seguir escribiendo
% «→» comodamente en el YAML.
\usepackage{newunicodechar}
\newunicodechar{→}{\ensuremath{\rightarrow}}
% Mismo problema con los simbolos de marca y vinetas: el «✓» de «una palomita ✓»
% salia como cuadrito vacio en el PDF de todas las guias que lo usaban.
\usepackage{pifont}
\newunicodechar{✓}{\ding{51}}
\newunicodechar{✗}{\ding{55}}
\newunicodechar{✔}{\ding{52}}
\newunicodechar{•}{\textbullet}
\newunicodechar{≥}{\ensuremath{\geq}}
\newunicodechar{≤}{\ensuremath{\leq}}

\setmainfont{Helvetica Neue}
\setsansfont{Helvetica Neue}
\setlength{\parindent}{0pt}
\setlength{\parskip}{6pt}
% Sin particion de palabras: en 6.º-7.º una palabra cortada por guion al final
% de linea («Comuni-cacion») frena la lectura mucho mas que una linea corta.
\hyphenpenalty=10000
\exhyphenpenalty=10000
\pretolerance=2000
\tolerance=3000
\setlength{\emergencystretch}{3em}
\linespread{1.10}
\renewcommand{\arraystretch}{1.25}
\setlist{leftmargin=5mm,itemsep=1mm,topsep=1mm}
\newcolumntype{Y}{>{\RaggedRight\arraybackslash}X}

\definecolor{milcMagenta}{HTML}{E6007E}
\definecolor{milcVino}{HTML}{5A0038}
\definecolor{milcTurquesa}{HTML}{0093A5}
\definecolor{milcVerde}{HTML}{58A923}
\definecolor{milcMostaza}{HTML}{E5B400}
\definecolor{milcOcre}{HTML}{B86600}
\definecolor{milcNegro}{HTML}{241816}
\definecolor{milcGris}{HTML}{F7F7F4}
\definecolor{milcLinea}{HTML}{E8E2DD}
\definecolor{milcGradePrimary}{HTML}{FF2D87}
\definecolor{milcGradeDark}{HTML}{9F1B56}
\definecolor{milcGradeSoft}{HTML}{FFE0EE}
\definecolor{milcGradeAccent}{HTML}{CFE7FF}
\definecolor{milcGradeText}{HTML}{FFFFFF}
\color{milcNegro}

\pagestyle{fancy}
\fancyhf{}
\lhead{\small Tecnología e Informática · Grado 10}
\rhead{\small Guía 28 · MILC v3}
\cfoot{\small\thepage}
\renewcommand{\headrulewidth}{0pt}

\newtcolorbox{titlebox}[1]{
  enhanced,colback=#1,colframe=#1,arc=7mm,boxrule=0pt,
  left=7mm,right=7mm,top=3mm,bottom=3mm,
  before skip=8pt,after skip=8pt,
  fontupper=\sffamily\bfseries\Large\color{white}
}

\newtcolorbox{softbox}[3]{
  enhanced,colback=#2,colframe=#1,borderline west={4pt}{0pt}{#1},
  arc=2mm,boxrule=.4pt,left=4mm,right=4mm,top=3mm,bottom=3mm,
  before skip=6pt,after skip=8pt,title={#3},coltitle=white,
  colbacktitle=#1,fonttitle=\sffamily\bfseries,fontupper=\RaggedRight,
  attach boxed title to top left={xshift=3mm,yshift=-2mm},
  boxed title style={arc=2mm, boxrule=0pt}
}

\newtcolorbox{drawbox}[2]{
  enhanced,colback=white,colframe=#1,arc=4mm,boxrule=1pt,
  left=5mm,right=5mm,top=4mm,bottom=4mm,before skip=8pt,after skip=8pt,
  title={#2},coltitle=white,colbacktitle=#1,fonttitle=\sffamily\bfseries,
  fontupper=\RaggedRight,
  attach boxed title to top left={xshift=4mm,yshift=-2mm},
  boxed title style={arc=3mm, boxrule=0pt}
}

\newtcolorbox{infoband}[1]{
  enhanced,colback=#1!8,colframe=#1!50,arc=2mm,boxrule=.5pt,
  left=4mm,right=4mm,top=2mm,bottom=2mm,before skip=4pt,after skip=4pt,
  fontupper=\small\RaggedRight
}

% Puente narrativo entre secciones (contrato editorial Conéctate).
% Da continuidad: "Acabas de X. Ahora vas a Y porque Z."
% Visualmente: banda izquierda en color del grado, texto en itálica pequeña.
\newtcolorbox{puentebox}{
  enhanced,colback=milcGradeSoft,colframe=milcGradeDark,
  borderline west={3pt}{0pt}{milcGradeDark},
  arc=0mm,boxrule=0pt,left=5mm,right=4mm,top=2.5mm,bottom=2.5mm,
  before skip=4pt,after skip=8pt,
  fontupper=\itshape\small\color{milcNegro}\RaggedRight
}
% La flecha va en modo matematico a proposito: Helvetica Neue NO tiene el glifo
% U+2192, asi que \textrightarrow se compilaba a NADA («Missing character: There
% is no → in font Helvetica Neue») y el puente salia sin flecha. En modo
% matematico la flecha viene de la fuente matematica y si se dibuja.
\newcommand{\puente}[1]{%
  \begin{puentebox}%
    \textcolor{milcGradeDark}{$\rightarrow$}\hspace{2mm}#1%
  \end{puentebox}%
}

\newcommand{\linead}[1]{\noindent\color{milcLinea}\rule{#1}{0.5pt}\color{milcNegro}}
\newcommand{\respuesta}[1]{\par\vspace{2mm}\foreach \n in {1,...,#1}{\noindent\color{milcLinea}\rule{\linewidth}{0.5pt}\par\vspace{5mm}}\color{milcNegro}}
\newcommand{\checkbox}{\textcolor{milcGradeDark}{\fbox{\phantom{x}}}\hspace{5pt}}

\newcommand{\phasecard}[4]{
\begin{tcolorbox}[enhanced,colback=#3,colframe=#2,arc=3mm,boxrule=.5pt,height=3.05cm,valign=center,halign=center,left=1.5mm,right=1.5mm,top=2mm,bottom=2mm]
{\sffamily\bfseries\fontsize{22}{24}\selectfont\textcolor{#2}{#1}}\par
{\sffamily\bfseries\small #4}
\end{tcolorbox}
}

% ── Piezas del rediseno 2026 (legibilidad 6.º-11.º) ───────────────────
% Banda de estacion: reemplaza el titlebox macizo. Titulo a la izquierda,
% dato practico (tiempo/modalidad) a la derecha, en una sola linea.
\newcommand{\milcestacion}[2]{%
\begin{tcolorbox}[enhanced,colback=#1,colframe=#1,arc=2mm,boxrule=0pt,
  left=5mm,right=5mm,top=2.6mm,bottom=2.6mm,before skip=11pt,after skip=7pt]
{\sffamily\bfseries\fontsize{15}{18}\selectfont\color{white}#2}
\end{tcolorbox}}

% Tarjeta de paso numerada: sustituye las tablas «Pilar / Aplicacion», que
% eran formato de consulta docente, no de estudiante. El numero va en un
% circulo sobre el borde para que la secuencia se lea de un vistazo.
\newcommand{\milcpaso}[3]{%
\begin{tcolorbox}[enhanced,colback=white,colframe=milcLinea,arc=1.5mm,boxrule=.9pt,
  left=14mm,right=4mm,top=3mm,bottom=3mm,before skip=4pt,after skip=4pt,
  fontupper=\RaggedRight,
  overlay={\node[anchor=north west,circle,fill=#2,text=white,inner sep=0pt,
    minimum size=7.4mm,font=\sffamily\bfseries\small]
    at ([xshift=3.6mm,yshift=-3.2mm]frame.north west) {#1};}]
#3
\end{tcolorbox}}

% Casilla marcable de tamano util con lapiz (la anterior era \fbox{\phantom{x}},
% de alto variable segun el texto que la seguia).
\newcommand{\milcmarca}{\framebox[3.8mm]{\rule{0pt}{3.4mm}}}

% Fila marcable de lista suelta (checks de la estacion 1).
\newcommand{\milccheckrow}[1]{%
\par\vspace{1.8mm}\noindent
\makebox[6.4mm][l]{\raisebox{-.55mm}{\textcolor{milcNegro}{\milcmarca}}}%
\parbox[t]{\dimexpr\linewidth-6.4mm\relax}{\RaggedRight #1}\par}

% Celda marcable dentro de la rubrica (tabla de autoevaluacion). Misma
% sangria francesa, pero pensada para vivir dentro de una columna X.
\newcommand{\milcopcion}[1]{%
\makebox[5.2mm][l]{\raisebox{-.55mm}{\textcolor{milcNegro}{\milcmarca}}}%
\parbox[t]{\dimexpr\linewidth-5.2mm\relax}{\RaggedRight #1}}

% ── Recursos visuales de la guía ──────────────────────────────────────
% Declarados en el bloque `recursos:` del YAML y emitidos por el builder.
% \guiaFigura   → imagen rasterizada o PDF (foto, carta celeste, montaje).
% \guiaDiagrama → fuente TikZ/circuitikz (.tex) incrustada como vector:
%                 es la vía para circuitos y esquemáticos editables.
% #1 = ruta relativa al .tex · #2 = pie (puede ir vacío)
\newcommand{\guiaFigura}[2]{%
\begin{center}
\includegraphics[width=0.88\linewidth,height=0.38\textheight,keepaspectratio]{#1}\\[1.5mm]
{\footnotesize\itshape\textcolor{milcNegro!75}{#2}}
\end{center}
}

\newcommand{\guiaDiagrama}[2]{%
\begin{center}
\input{#1}\\[1.5mm]
{\footnotesize\itshape\textcolor{milcNegro!75}{#2}}
\end{center}
}

\begin{document}
\thispagestyle{empty}

% ── Portada ───────────────────────────────────────────────────────────
% Antes: un rectangulo de color a pagina completa. Se veia bien en pantalla,
% pero en la fotocopiadora del colegio se comia el toner y salia gris sucio.
% Ahora la banda ocupa el tercio superior y el resto va en blanco.
\begin{tikzpicture}[remember picture,overlay]
  \path[fill=milcGradePrimary]
    (current page.north west) rectangle ([yshift=-10.6cm]current page.north east);

  \node[anchor=north west,text=milcGradeText] at ([xshift=2.0cm,yshift=-1.55cm]current page.north west)
    {\sffamily\bfseries\fontsize{12}{14}\selectfont GRADO};
  \node[anchor=north west,text=milcGradeText,opacity=.85] at ([xshift=2.0cm,yshift=-2.35cm]current page.north west)
    {\sffamily\bfseries\fontsize{8.5}{10}\selectfont SERIE GUÍAS · TECNOLOGÍA E INFORMÁTICA};
  \node[anchor=north east,text=milcGradeText,opacity=.85,align=right] at ([xshift=-2.0cm,yshift=-1.55cm]current page.north east)
    {\sffamily\bfseries\fontsize{9}{12}\selectfont I.E. Sor María Juliana\\Cartago, Valle del Cauca};

  \node[anchor=west,text=milcGradeText] (gradonum) at ([xshift=1.85cm,yshift=-7.1cm]current page.north west)
    {\sffamily\bfseries\fontsize{112}{112}\selectfont 10};
  % El circulito de grado (°) solo aplica a las guias de grado («11°»); en
  % Bebras y Semillero el numero grande es un momento/modulo, no un grado.
  % Se posiciona relativo al nodo `gradonum`, no en coordenadas absolutas.
  \draw[milcGradeText,line width=1.6pt]
    ([xshift=2.0mm,yshift=-4.5mm]gradonum.north east) circle (.15cm);

  \node[anchor=west,text=milcGradeText,text width=11.4cm,align=left] at ([xshift=1.5cm,yshift=0.5cm]gradonum.east)
    {
     {\sffamily\bfseries\fontsize{9}{11}\selectfont\MakeUppercase{Período 3 · Ofimática con IA: contabilidad, datos y emprendimiento}}\\[3mm]
     {\sffamily\bfseries\fontsize{21}{25}\selectfont\setlength{\baselineskip}{27pt}Comunicación\\[4pt]empresarial profesional\\[4pt]con IA}\\[3.5mm]
     {\sffamily\bfseries\fontsize{15}{18}\selectfont Guía 28-10-TIC}
    };
\end{tikzpicture}

\vspace*{9.4cm}
\vfill

{\sffamily\bfseries\fontsize{8}{10}\selectfont\textcolor{milcGradeDark}{LO QUE TE LLEVAS HOY}}

\begin{tcolorbox}[enhanced,colback=white,colframe=milcNegro,arc=2mm,boxrule=1.6pt,
  left=5mm,right=5mm,top=4mm,bottom=4mm,before skip=3pt,after skip=12pt,
  fontupper=\RaggedRight\fontsize{11}{15.5}\selectfont]
3 piezas de comunicación empresarial del microemprendimiento producidas con asistencia de IA (a) propuesta comercial breve de 1 página dirigida a cliente potencial real o hipotético, (b) correo de presentación profesional para abrir relación comercial con cliente, (c) infografía o publicación visual para redes sociales con el resumen de la propuesta de valor + declaración explícita por cada pieza de qué generó la IA, qué editaste tú y por qué
\end{tcolorbox}

\begin{tcolorbox}[enhanced,colback=milcGris,colframe=milcGris,arc=2mm,boxrule=0pt,
  left=5mm,right=5mm,top=3.5mm,bottom=3.5mm,after skip=12pt]
{\sffamily\bfseries\fontsize{8}{10}\selectfont\textcolor{milcNegro!55}{ESTA GUÍA ES DE}}\\[3mm]
\begin{tabularx}{\linewidth}{X p{3.2cm} p{3.2cm}}
\linead{\linewidth} & \linead{3.0cm} & \linead{3.0cm}\\[-1mm]
{\fontsize{8.5}{10}\selectfont\textcolor{milcNegro!55}{Nombre completo}} &
{\fontsize{8.5}{10}\selectfont\textcolor{milcNegro!55}{Grupo}} &
{\fontsize{8.5}{10}\selectfont\textcolor{milcNegro!55}{Fecha}}\\
\end{tabularx}
\end{tcolorbox}

\vfill

\noindent\textcolor{milcLinea}{\rule{\linewidth}{1pt}}\\[2mm]
{\sffamily\bfseries\fontsize{9.5}{12}\selectfont Autor: Dr. Álvaro Cárdenas Orozco}\\[1mm]
{\sffamily\fontsize{8.5}{11}\selectfont\textcolor{milcNegro!55}{Metodología MILC · apertura ancestral · pensamiento computacional · triángulo Dussel–estoicismo–Floridi}}

\newpage

\section*{Apertura: Comunicación empresarial profesional con IA}

\begin{softbox}{milcGradeDark}{milcGradeSoft}{Saber ancestral}
En las tiendas, panaderías y modisterías del centro de Cartago, había un concepto que sostenía el comercio del barrio durante décadas: \textbf{la palabra empeñada}. Cuando un tendero le decía a un cliente regular \emph{``el lunes le tengo la harina especial''}, esa frase no era promesa vaga: era \textbf{compromiso público}. Si no llegaba el lunes, el tendero perdía credibilidad y el cliente cambiaba de tienda. La palabra empeñada tenía 3 propiedades inquebrantables: (1) \textbf{Claridad}: nada de ambigüedades. Día, hora, producto específicos. (2) \textbf{Realismo}: solo se prometía lo que se podía cumplir. (3) \textbf{Honor}: el cumplimiento era cuestión de respeto, no de contrato legal. Los clientes regulares del barrio no firmaban papeles; trabajaban con palabra empeñada porque sabían que el tendero respondería con su nombre. Cuando alguien decía \emph{``ese tendero tiene palabra''}, era el mayor elogio comercial posible. Esa misma palabra empeñada es la base de la comunicación empresarial profesional moderna: cada propuesta, cada correo, cada publicación son \textbf{compromisos públicos del emprendedor}. La IA ayuda a redactarlos; tú asumes lo prometido.
\end{softbox}

\begin{softbox}{milcTurquesa}{milcGris}{Saber contemporáneo}
La \textbf{comunicación empresarial profesional} es el conjunto de piezas escritas y visuales que un emprendedor produce para relacionarse con clientes, aliados, inversionistas. Sus 3 piezas más usadas en microemprendimiento son: (1) \textbf{Propuesta comercial breve}: 1 página dirigida a cliente potencial con problema + solución + propuesta + precio + llamada a la acción. (2) \textbf{Correo de presentación}: contacto inicial profesional con asunto claro, presentación breve, propuesta de valor, pedido de reunión o respuesta. (3) \textbf{Publicación visual para redes}: infografía o post con la propuesta de valor en formato consumible (visual, claro, llamado a acción). La regla profesional: \textbf{las 3 piezas deben ser coherentes entre sí}: si tu correo dice X y tu propuesta dice Y, el cliente se confunde. La \textbf{IA como asistente} acelera la redacción de las 3, pero el emprendedor: (a) ajusta tono al contexto colombiano, (b) garantiza coherencia con el Canvas, (c) asume compromisos prometidos. Herramientas: ChatGPT/Gemini para texto, Canva para visual.
\end{softbox}

\begin{softbox}{milcMostaza}{milcGris}{Pregunta-puente}
\textit{¿Qué sabía el tendero al cumplir su palabra empeñada el lunes, que el novato olvida cuando promete en correo lo que no podrá cumplir? ¿Y por qué la coherencia entre las 3 piezas es lo que diferencia al emprendedor profesional del improvisado?}
\end{softbox}

\begin{softbox}{milcVerde}{milcGris}{Saber y saber hacer}
Al terminar podrás: (1) \textbf{aplicar} las habilidades de S8 del P2 (correo profesional con IA) al contexto empresarial; (2) \textbf{analizar} las 3 piezas de comunicación buscando coherencia y honestidad; (3) \textbf{crear} las 3 piezas con asistencia de IA y voz propia; (4) \textbf{evaluar} si las 3 piezas funcionan juntas como sistema coherente.
\end{softbox}



\small
\begin{tabularx}{\textwidth}{>{\bfseries}p{3.0cm}Y}
\rowcolor{milcGradePrimary!12}Autor & Dr. Álvaro Cárdenas Orozco\\
DBA & Aplicación de ofimática asistida por IA a contabilidad básica y emprendimiento (MEN, T\&I)\\
Referentes & Lean Startup · Phronesis financiera · Ética del dato empresarial\\
Producto final & 3 piezas de comunicación empresarial del microemprendimiento producidas con asistencia de IA (a) propuesta comercial breve de 1 página dirigida a cliente potencial real o hipotético, (b) correo de presentación profesional para abrir relación comercial con cliente, (c) infografía o publicación visual para redes sociales con el resumen de la propuesta de valor + declaración explícita por cada pieza de qué generó la IA, qué editaste tú y por qué\\
\end{tabularx}
\normalsize

\puente{Antes de empezar, mira el viaje completo. En la sesión 7 cerraste el Canvas. Hoy creas las \textbf{piezas comunicativas} que llevan el Canvas al cliente. La sesión 9 será el proyecto integrador completo y la 10 sustentación.}

\milcestacion{milcGradeDark}{Ruta MILC: así vamos a aprender}
\begin{tabularx}{\textwidth}{>{\centering\arraybackslash}X>{\centering\arraybackslash}X>{\centering\arraybackslash}X>{\centering\arraybackslash}X}
\phasecard{1}{milcVerde}{milcGris}{Escucha\\[-1mm]\small Observo, escucho y pregunto.} &
\phasecard{2}{milcTurquesa}{milcGris}{Entender\\[-1mm]\small Sistematizo y ordeno.} &
\phasecard{3}{milcMagenta}{milcGris}{Hacer\\[-1mm]\small Construyo y aplico.} &
\phasecard{4}{milcVino}{milcGris}{Cierre\\[-1mm]\small Evalúo y reflexiono.}
\end{tabularx}


\puente{Antes de teorizar, vas a hacer una \emph{revisión rápida}: lee tu Canvas de la sesión 7 e identifica la propuesta de valor en 1 frase. Esa frase es el corazón de las 3 piezas que vas a crear hoy.}

\milcestacion{milcVerde}{Estación 1 de 3 · Escucha}

\begin{softbox}{milcVerde}{milcGris}{Escena para escuchar}
\textbf{Actividad 1 · APLICA --- Propuesta de valor en 1 frase} (15 min · individual). Toma tu Canvas de la sesión 7 e identifica la propuesta de valor de tu microemprendimiento. Reduce a 1 frase específica de máximo 15-20 palabras. Estructura: \emph{``Ofrezco [producto/servicio] a [cliente específico] que resuelve [problema concreto] mediante [diferencial clave]''}. Ejemplo: \emph{``Ofrezco empanadas vegetarianas a estudiantes universitarios del Valle que quieren almuerzo rápido sin carne mediante delivery por WhatsApp en 30 minutos''}.
\end{softbox}

{\sffamily\bfseries\fontsize{8}{10}\selectfont\textcolor{milcNegro!55}{MARCA CUANDO LO HAYAS HECHO}}
\milccheckrow{Tomé propuesta de valor del Canvas.}
\milccheckrow{La reduje a 1 frase de máximo 20 palabras.}
\milccheckrow{La frase incluye producto, cliente, problema, diferencial.}

\begin{infoband}{milcVerde}


\textbf{Lo que va al cuaderno (Actividad 1 · APLICA).} Título: «Actividad 1 --- Propuesta en 1 frase». \textbf{Formato:} la frase final + 2-3 versiones previas descartadas con razón. \textbf{Sabes que terminaste cuando} esa frase suena clara y específica.
\end{infoband}

\puente{Ya tienes propuesta de valor en 1 frase. Ahora vas a entender la \textbf{anatomía profesional} de las 3 piezas de comunicación empresarial y cómo mantener coherencia entre ellas.}

\milcestacion{milcTurquesa}{Estación 2 de 3 · Entender}

\begin{softbox}{milcTurquesa}{milcGris}{Lo que vas a entender}
Tu propuesta en 1 frase es el corazón. Ahora necesitas la \textbf{anatomía profesional} de las 3 piezas de comunicación empresarial.
\end{softbox}

\milcpaso{1}{milcTurquesa}{Las \textbf{3 piezas de comunicación empresarial} se descomponen así: (1) \textbf{Propuesta comercial breve} (1 página): título atractivo, problema del cliente en 2-3 líneas, solución que ofreces en 2-3 líneas, beneficios principales en 3-4 puntos, precio claro, llamada a la acción específica. Tono: cálido pero profesional. (2) \textbf{Correo de presentación}: asunto claro y específico, saludo apropiado, 3-4 párrafos cortos (quién eres + propuesta de valor + pedido), cierre con llamada a acción concreta. Tono: profesional pero accesible. (3) \textbf{Publicación visual para redes}: imagen o infografía con propuesta de valor en gran formato, beneficios principales como 3-4 íconos o frases cortas, llamada a acción visible. Diseñada en Canva con plantillas gratuitas.}
\milcpaso{2}{milcTurquesa}{La \textbf{coherencia entre las 3 piezas} es crucial. Las 3 deben: (a) comunicar la misma propuesta de valor con palabras similares (no idénticas, pero alineadas). (b) usar el mismo tono general. (c) hacer la misma promesa al cliente. (d) llamar a la misma acción final. Si la propuesta dice \emph{``delivery en 30 minutos''}, el correo y la publicación también deben decirlo. La incoherencia entre piezas confunde al cliente y erosiona la palabra empeñada.}
\milcpaso{3}{milcTurquesa}{El \textbf{flujo con IA} para cada pieza sigue 4 pasos: (1) escribir prompt completo con las 5 partes (rol, contexto, tarea, formato, restricciones), incluyendo tu propuesta de valor en 1 frase. (2) recibir borrador de IA en ChatGPT/Gemini/Bing. (3) editar a mano: ajustar tono colombiano, eliminar clichés, garantizar coherencia con las otras piezas. (4) verificar que la pieza compromete lo que efectivamente puedes cumplir. La phronesis del emprendedor consiste en \textbf{no usar IA para prometer lo que no podrás entregar}.}
\milcpaso{4}{milcTurquesa}{Algoritmo: (1) revisar propuesta de valor en 1 frase. (2) escribir prompt para propuesta comercial. (3) recibir borrador IA, editar, finalizar. (4) escribir prompt para correo de presentación. (5) recibir borrador, editar, finalizar. (6) escribir prompt para concepto de publicación visual. (7) usar Canva con plantilla para crear publicación. (8) revisar coherencia entre las 3 piezas.}

\begin{drawbox}{milcTurquesa}{Anatomía de las 3 piezas}
\textbf{Propuesta comercial:} 1 pp con problema-solución-precio. \\[1mm] \textbf{Correo de presentación:} asunto + 3-4 párrafos + llamada. \\[1mm] \textbf{Publicación redes:} infografía Canva con propuesta visual. \\[1mm] \textbf{Coherencia:} misma propuesta, mismo tono, misma promesa.
\end{drawbox}

\begin{softbox}{milcMostaza}{milcGris}{Errores comunes y su corrección}
(1) \textbf{Incoherencia entre piezas}: la propuesta dice algo, el correo otro. (2) \textbf{Prometer lo que no se puede cumplir}: la IA puede sugerir promesas atractivas pero irrealistas. (3) \textbf{Tono inconsistente}: propuesta formal con correo casual. (4) \textbf{Sin llamada a acción clara}: la pieza termina sin pedir nada concreto. (5) \textbf{Publicación visual sin propuesta clara}: bonito pero vacío.
\end{softbox}

\begin{infoband}{milcTurquesa}


\textbf{Lo que va al cuaderno (Actividad 2 · ANALIZA).} Título: «Actividad 2 --- Anatomía comunicación empresarial». \textbf{Formato:} (a) 3 piezas con estructura; (b) regla de coherencia; (c) los 5 errores con marca. \textbf{Sabes que terminaste cuando} sabes producir las 3 piezas sin abrir esta guía.
\end{infoband}

\puente{Tienes el método. Toca aplicarlo: creas las 3 piezas con asistencia de IA, garantizas coherencia entre ellas, declaras uso de IA en cada una.}

\milcestacion{milcMagenta}{Estación 3 de 3 · Hacer}

\begin{softbox}{milcMagenta}{milcGris}{Cómo vas a construir tu producto}
\textbf{Actividad 3 · CREA + EVALÚA --- 3 piezas + declaración de IA} (40 min · individual con dispositivo). Hora de aplicar. Creas las 3 piezas con asistencia de IA, garantizas coherencia, declaras uso de IA en cada una.
\end{softbox}

\milcpaso{1}{milcMagenta}{Tus 3 piezas se construyen en 6 pasos: (1) escribir prompt completo para propuesta comercial: \emph{``Actúa como mentor de emprendimiento. Mi propuesta de valor es [frase de S1]. Escribe propuesta comercial de 1 página dirigida a [cliente potencial] con problema, solución, 3 beneficios, precio y llamada a acción''}. (2) editar el borrador IA. (3) prompt para correo de presentación. (4) editar. (5) prompt para concepto de publicación visual + crear en Canva con plantilla. (6) revisar coherencia entre las 3.}
\milcpaso{2}{milcMagenta}{Sigue el patrón \textbf{``palabra empeñada digital''}: cada palabra de las 3 piezas es compromiso público. Si tu propuesta promete \emph{``entrega en 30 minutos''}, mejor que puedas cumplirlo. La IA puede sugerir promesas atractivas pero irrealistas; tu trabajo como emprendedor es \textbf{rechazar lo que no podrías cumplir}. La phronesis comercial consiste en \textbf{prometer poco y cumplir mucho}.}
\milcpaso{3}{milcMagenta}{La \textbf{declaración de uso de IA por pieza} tiene 3 elementos: (a) qué modelo IA usaste (ChatGPT, Gemini, etc.). (b) qué porcentaje del texto generó la IA y qué porcentaje editaste tú. (c) qué partes específicas cambiaste y por qué. Esa declaración por pieza es transparencia profesional contemporánea.}
\milcpaso{4}{milcMagenta}{Algoritmo: (1) propuesta comercial. (2) correo presentación. (3) publicación visual. (4) verificación de coherencia. (5) declaración IA por cada pieza.}

\begin{drawbox}{milcMagenta}{Checklist antes de entregar}
\checkbox Propuesta comercial de 1 página completa. \\[1mm] \checkbox Correo de presentación con asunto y llamada. \\[1mm] \checkbox Publicación visual en Canva. \\[1mm] \checkbox 3 piezas coherentes entre sí. \\[1mm] \checkbox Promesas realistas en cada pieza. \\[1mm] \checkbox Declaración de IA por cada pieza.
\end{drawbox}

\begin{softbox}{milcMagenta}{milcGris}{Plantilla de trabajo}
\textbf{Pieza 1.} \textit{Propuesta comercial 1 página.} \\[1mm] \textbf{Pieza 2.} \textit{Correo de presentación.} \\[1mm] \textbf{Pieza 3.} \textit{Publicación visual.} \\[1mm] \textbf{Coherencia.} \textit{Misma propuesta + tono + promesa.} \\[1mm] \textbf{IA.} \textit{Declaración por cada pieza.}
\end{softbox}

\begin{infoband}{milcMagenta}


\textbf{Lo que va al cuaderno (Actividad 3 · CREA + EVALÚA).} Título: «Actividad 3 --- 3 piezas empresariales». \textbf{Formato:} (a) las 3 piezas; (b) verificación coherencia; (c) declaración IA. \textbf{Sabes que terminaste cuando} las 3 piezas podrían enviarse a cliente potencial real.
\end{infoband}

\puente{Lo que sigue resume \emph{qué} entregas hoy: 3 piezas coherentes + declaración de IA. El criterio principal: que las 3 piezas podrían enviarse hoy mismo a un cliente potencial real.}

\milcestacion{milcMostaza}{Producto de la sesión}
\textbf{3 piezas empresariales coherentes + declaración de IA}.
\begin{softbox}{milcMostaza}{milcGris}{Criterios de calidad}
(1) Propuesta comercial 1 página. (2) Correo de presentación profesional. (3) Publicación visual en Canva. (4) Coherencia entre las 3 piezas. (5) Promesas realistas. (6) Declaración IA por pieza.
\end{softbox}

\puente{Antes de cerrar, mira la comunicación desde las cinco dimensiones humanas. Mantener palabra empeñada digital es respeto por el cliente; prometer lo que no se puede cumplir es traición al oficio.}

\milcestacion{milcVino}{Evaluación liberadora · cinco dimensiones}

\begin{softbox}{milcVino}{milcGris}{Sentido de la evaluación}
Evaluar no es solo revisar una nota. Es reconocer qué aprendí, qué cambió en mí, qué emociones manejé, cómo me relaciono con la ciudad y la comunidad, y qué le dejaré a quien viene detrás.
\end{softbox}

\textbf{Mi reflexión liberadora en cinco dimensiones.} Completa cada frase iniciada:

\small
\begin{tabularx}{\textwidth}{>{\bfseries\RaggedRight\arraybackslash}p{3.4cm}Y>{\RaggedRight\arraybackslash}p{4.6cm}}
\rowcolor{milcVino}\color{white}Dimensión & \color{white}\bfseries Mi respuesta (frame starter) & \color{white}\bfseries Algo que descubrí\\
1. Personal & Lo que más cambió en mí fue \linead{6.5cm} & \linead{4.2cm}\\
2. Emocional & La emoción más fuerte fue \linead{2.5cm} y la regulé \linead{3cm} & \linead{4.2cm}\\
3. Ciudadana & En democracia esto importa porque \linead{6cm} & \linead{4.2cm}\\
4. Local (Cartago) & En mi barrio sería útil para \linead{6.5cm} & \linead{4.2cm}\\
5. Intergeneracional & A mi abuela/abuelo le diría \linead{6.5cm} & \linead{4.2cm}\\
\end{tabularx}
\normalsize

\begin{infoband}{milcVino}
\textbf{Para la próxima sustentación.} Vuelve a esta tabla en seis meses. Verás que las respuestas cambian — no porque lo aprendido cambie, sino porque tú cambias. La evaluación liberadora es una conversación contigo a través del tiempo.
\end{infoband}


\puente{Y para terminar, tres voces sobre la palabra empeñada digital. Dussel sobre la comunicación que dignifica al cliente, Marco Aurelio sobre la coherencia entre las piezas, Floridi sobre la transparencia empresarial como ética emprendedora.}

\milcestacion{milcGradeDark}{Triángulo de pensamiento · cierre filosófico}

\begin{softbox}{milcVino}{milcGris}{Dussel · lente del nosotros}
\textit{La palabra empeñada digital es respeto por el cliente; prometer lo no cumplible es violencia comercial moderna.}

Cuando tus piezas prometen al cliente algo que no podrás cumplir, traicionas la palabra empeñada. La filosofía de la liberación pide lo opuesto: \textbf{prometer solo lo cumplible}. Las microempresas del barrio que duran décadas son las que mantienen palabra; las que truenan son las que prometen lo que no entregan. La phronesis empresarial consiste en \textbf{cumplir más de lo prometido, no menos}.

\textbf{¿Mis 3 piezas prometen lo que efectivamente puedo cumplir, o prometen más por verse atractivas?}
\end{softbox}
\linead{\linewidth}\\[1mm]
\linead{\linewidth}

\begin{softbox}{milcMostaza}{milcGris}{Estoicismo · Marco Aurelio · cuidado interior}
\textit{La coherencia entre las piezas comunicativas es virtud profesional; la incoherencia es debilidad disfrazada de creatividad.}

Es tentador hacer cada pieza con tono distinto creyendo que \emph{``así llega a más audiencias''}. La sabiduría estoica recomienda lo opuesto: \emph{coherencia es respeto por el cliente que ve todas las piezas}. La phronesis del oficio consiste en \textbf{decir lo mismo con palabras adaptadas, no cosas distintas según conveniencia}.

\textbf{¿Mis 3 piezas son coherentes, o cada una promete cosa distinta?}
\end{softbox}
\linead{\linewidth}\\[1mm]
\linead{\linewidth}

\begin{softbox}{milcTurquesa}{milcGris}{Floridi · lente de la infoesfera}
\textit{La transparencia empresarial con declaración honesta de IA es ética emprendedora del siglo XXI.}

En la era de la IA generativa, las microempresas que declaran abiertamente uso de IA se posicionan como ético emergente. Ocultar uso de IA mientras se promete \emph{``hecho a mano''} es fraude moderno. Tu declaración por pieza es práctica que rige todo emprendimiento serio contemporáneo.

\textbf{¿Mi comunicación empresarial declara honestamente uso de IA, o esconde la asistencia?}
\end{softbox}
\linead{\linewidth}\\[1mm]
\linead{\linewidth}

\begin{infoband}{milcNegro}
\textbf{Nota para el docente.} Las tres formulaciones del triángulo son la síntesis del autor de la guía, no citas textuales de los pensadores. Si vas a citarlos en clase, remite a la obra original.
\end{infoband}

\milcestacion{milcVino}{Mi autoevaluación · marca una casilla por fila}

{\small
\setlength{\extrarowheight}{2.2pt}
\begin{tabularx}{\textwidth}{>{\RaggedRight\arraybackslash\bfseries}p{3.0cm}YYY}
\rowcolor{milcVino}\color{white}\bfseries Criterio & \color{white}\bfseries Lo logré & \color{white}\bfseries En proceso & \color{white}\bfseries Necesito apoyo\\[.6mm]
Comprensión del tema & \milcopcion{Explico las ideas con mis palabras.} & \milcopcion{Reconozco las ideas, falta precisión.} & \milcopcion{Necesito repasar lo básico.}\\[.8mm]
\rowcolor{milcGris}Pensamiento computacional & \milcopcion{Usé los pilares al armar mi producto.} & \milcopcion{Usé algunos pilares.} & \milcopcion{Necesito releer la estación 2.}\\[.8mm]
Aplicación y producto & \milcopcion{Mi producto cumple los criterios.} & \milcopcion{Está armado, con ajustes pendientes.} & \milcopcion{Aún está en construcción.}\\[.8mm]
\rowcolor{milcGris}Reflexión 5 dimensiones & \milcopcion{Respondí cada una con honestidad.} & \milcopcion{Respondí algunas con detalle.} & \milcopcion{Mis respuestas son muy generales.}\\[.8mm]
Triángulo de pensamiento & \milcopcion{Conecté las 3 voces con mi trabajo.} & \milcopcion{Conecté al menos una.} & \milcopcion{No las relaciono con mi proceso.}\\[.8mm]
\end{tabularx}}

\vspace{3mm}
\textbf{Hoy entendí que:}\\[1mm]
\linead{\linewidth}\\[1mm]
\linead{\linewidth}

\textbf{Mi compromiso para la próxima sesión es:}\\[1mm]
\linead{\linewidth}\\[1mm]
\linead{\linewidth}

\begin{softbox}{milcMostaza}{milcGris}{Cita de despedida · Marco Aurelio}
\textit{``El obstáculo en el camino se vuelve el camino. Lo que estorba al camino se vuelve camino.''}\\[1mm]
Cada error de hoy, bien observado, se convierte en pieza del camino que pisarás mañana.
\end{softbox}

\small
\begin{tabularx}{\textwidth}{>{\bfseries}p{3.6cm}Y}
\rowcolor{milcGradeDark!12}Nombre completo: & \linead{10cm}\\
Fecha: & \linead{4cm} \quad Curso/grupo: \linead{4cm}\\
Mi firma: & \linead{9cm}\\
\end{tabularx}
\normalsize

\begin{infoband}{milcGradeDark}
\textbf{Para la próxima sesión.} Llega con tus 3 piezas coherentes y declaración de IA. En la sesión 9 vas a integrar todo el periodo en un \textbf{proyecto integrador completo del micro-emprendimiento}: Canvas + flujo de caja + estudio mercado + comunicación + plan de validación. La comunicación de hoy es una pieza; el proyecto integrador es la integración total.
\end{infoband}

\end{document}
